\section*{Chapitre \ref{ch:pl:messaging-and-event-driven-architecture} --- Messagerie et architecture pilotée par les événements}

\begin{solution}{exo:pl:messaging-and-event-driven-architecture:1}
L'ordre n'est conservé qu'au sein d'une partition. Un regroupement par compte place toutes les exécutions d'un compte dans une même partition, de sorte que le \omterm{def:pl:position-and-pnl-services:service}{service de positions} les applique
dans l'ordre où elles se sont produites~; une distribution circulaire les disperserait sur des partitions lues à des vitesses différentes.
\end{solution}

\begin{solution}{exo:pl:messaging-and-event-driven-architecture:2}
Le consommateur reprend à la dernière \omterm{def:pl:messaging-and-event-driven-architecture:log}{position du consommateur} validée, soit au début du lot, et applique les 100 messages~: les 30 déjà appliqués sont donc traités deux fois.
\end{solution}

\begin{solution}{exo:pl:messaging-and-event-driven-architecture:3}
Distribution circulaire~: trois, trois et deux partitions (0, 3, 6~; 1, 4, 7~; 2, 5). Un quatrième membre déclenche un rééquilibrage à deux par membre~; le nouveau
propriétaire de chaque partition reprend à la \omterm{def:pl:messaging-and-event-driven-architecture:log}{position du consommateur} validée du groupe, si bien que tout ce que l'ancien propriétaire a traité sans la valider est traité de nouveau.
\end{solution}

\begin{solution}{exo:pl:messaging-and-event-driven-architecture:4}
Un plantage survient uniformément dans un lot de 100, en moyenne à mi-parcours. Valider d'abord a déjà fait avancer la position au-delà du lot, si bien que la
moitié non traitée est perdue~; valider ensuite ne l'a pas fait, si bien que la moitié traitée est répétée. Les 50\,323 et 48\,652 observés pour 1\,000 plantages
correspondent à ce demi-lot.
\end{solution}

\begin{solution}{exo:pl:messaging-and-event-driven-architecture:5}
Dans l'instantané, écrit de façon atomique avec l'état qu'il décrit~: au redémarrage, on charge l'instantané et on reprend à la \omterm{def:pl:messaging-and-event-driven-architecture:log}{position du consommateur} qu'il contient,
en rejouant tout ce qui suit. Une position validée ailleurs ne correspondrait pas à l'état.
\end{solution}

\begin{solution}{exo:pl:messaging-and-event-driven-architecture:6}
Un relais qui publie un événement puis tombe avant d'avoir marqué la ligne comme envoyée le republie après le redémarrage. Le producteur idempotent (le journal
ignore un numéro d'outbox déjà connu) ou le consommateur (qui déduplique sur la clé de l'événement) l'élimine.
\end{solution}

\begin{solution}{exo:pl:messaging-and-event-driven-architecture:7}
Conserver l'ensemble des \omterm{def:pl:position-and-pnl-services:execid}{identifiants d'exécution} appliqués avec les positions et ignorer tout identifiant déjà présent~; avec l'ensemble et les positions mis à jour
ensemble, aucune exécution n'est perdue (la \omterm{def:pl:messaging-and-event-driven-architecture:log}{position du consommateur} est toujours validée après) et aucune ne compte deux fois~: zéro et zéro, comme dans le mode atomique.
\end{solution}

\begin{solution}{exo:pl:messaging-and-event-driven-architecture:8}
La garantie de \omterm{def:pl:messaging-and-event-driven-architecture:exactly}{traitement exactement une fois} d'un produit de messagerie couvre ce qui se passe à l'intérieur de celui-ci --- écritures dédupliquées dans son journal, transactions entre ses partitions ---
et non l'effet d'un message sur un autre système. La base de positions en est un autre~: ses mises à jour et la \omterm{def:pl:messaging-and-event-driven-architecture:log}{position du consommateur} doivent être rendues
atomiques, ou son traitement rendu idempotent, par le consommateur.
\end{solution}

\begin{solution}{pb:pl:messaging-and-event-driven-architecture:1}
\begin{enumerate}
\item Connaissance, temps et débit~: les éditeurs ne connaissent pas les consommateurs, les consommateurs peuvent être arrêtés puis rattraper leur retard, et un consommateur lent
ne ralentit pas l'éditeur. Le coût est un graphe d'appels visible, et les pannes se déplacent vers la livraison.
\item Un courtier lorsque le stockage, le rejeu et l'évolution des consommateurs comptent~; sans courtier (multicast, mémoire partagée, une bibliothèque comme Aeron) lorsque la
latence compte et que les extrémités sont connues.
\item Par son offset dans sa partition~; la partition résulte d'un hachage stable de la clé.
\item Les partitions d'un topic, un membre par partition, et une \omterm{def:pl:messaging-and-event-driven-architecture:log}{position du consommateur} validée par partition~; un rééquilibrage réattribue les partitions lorsque les membres changent.
\item Des segments entiers d'anciens messages~; un nouvel abonné qui en a besoin doit partir d'un instantané qui enregistre l'offset couvert.
\item Le traitement et la validation de la \omterm{def:pl:messaging-and-event-driven-architecture:log}{position du consommateur}~; un plantage peut survenir pendant le traitement ou entre les deux.
\item 967 dans le mode validation d'abord, 1\,000 dans les autres~: valider d'abord ne répète jamais de messages, si bien qu'une journée compte moins d'étapes et que certains
points de plantage prévus tombent après sa fin.
\item Valider d'abord~: 50\,323 perdus, aucun dupliqué~; valider ensuite~: aucun perdu, 48\,652 dupliqués~; atomique~: aucun des deux.
\item Jusqu'à 5\,100 actions (médiane 3\,500) en validant d'abord, jusqu'à 6\,500 (médiane 3\,800) en validant ensuite, zéro en atomique.
\item Les positions et la \omterm{def:pl:messaging-and-event-driven-architecture:log}{position du consommateur} suivante sont écrites dans une même transaction~; un plantage annule les deux, et le consommateur reprend là où l'état
s'arrête.
\item Ce sont deux systèmes, la base de données et le journal, sans aucune transaction qui les couvre tous les deux.
\item Écrire puis publier~: 22 événements perdus~; publier puis écrire~: 22 événements fantômes.
\item Des répétitions (22 sans clés d'idempotence), éliminées par le producteur idempotent du journal ou par le consommateur.
\item L'état est stocké comme la suite des événements qui l'ont modifié~; le journal doit conserver chaque événement, ou un instantané accompagné de son décalage.
\item Les messages qui suivent un message toxique dans sa partition.
\item \textbf{Résultat nommé.} Pour 1\,000 plantages, valider la \omterm{def:pl:messaging-and-event-driven-architecture:log}{position du consommateur} avant le traitement (au plus une fois) perd 50\,323 exécutions et n'en duplique aucune~;
valider ensuite (au moins une fois) n'en perd aucune et en duplique 48\,652~; stocker positions et \omterm{def:pl:messaging-and-event-driven-architecture:log}{position du consommateur} dans une même transaction ne perd ni ne duplique
rien. L'erreur de position de fin de journée la pire sur cent jours est de 5\,100, 6\,500 et 0 actions.
\item Le mode atomique (ou idempotent) pour les positions~; un tableau de bord peut valider d'abord, car il n'a besoin que du dernier état.
\item Son propre journal~: écritures dédupliquées et transactions entre ses partitions, et non la base de données du consommateur.
\item Rapprocher en continu les positions avec la drop copy, par compte et par action, et alerter à la moindre différence.
\item Du côté du consommateur, qui rend son effet idempotent ou atomique avec sa \omterm{def:pl:messaging-and-event-driven-architecture:log}{position du consommateur}.
\end{enumerate}
\end{solution}

\begin{solution}{iq:pl:messaging-and-event-driven-architecture:1}
Au plus une fois~: la position est enregistrée avant le traitement, si bien qu'un plantage perd des messages. Au moins une fois~: enregistrée après, si bien qu'un plantage les répète.
\omterm{def:pl:messaging-and-event-driven-architecture:exactly}{Traitement exactement une fois}~: un effet, non une propriété du transport, obtenu par un traitement idempotent ou par un stockage atomique de l'état et de la position.
\iqlookfor{Où la \omterm{def:pl:messaging-and-event-driven-architecture:log}{position du consommateur} est enregistrée.}
\end{solution}

\begin{solution}{iq:pl:messaging-and-event-driven-architecture:2}
Valider la \omterm{def:pl:messaging-and-event-driven-architecture:log}{position du consommateur} dans la même transaction de base de données que les positions, ou stocker avec elles les \omterm{def:pl:position-and-pnl-services:execid}{identifiants d'exécution} appliqués et ignorer les répétitions~; ne jamais se
fier au seul transport.
\iqlookfor{\omterm{def:pl:messaging-and-event-driven-architecture:log}{Position du consommateur} atomique ou idempotence.}
\end{solution}

\begin{solution}{iq:pl:messaging-and-event-driven-architecture:3}
Une \omterm{def:pl:messaging-and-event-driven-architecture:outbox}{boîte d'envoi transactionnelle}~: l'événement est inséré dans la même transaction que la ligne~; un relais publie les événements non envoyés dans l'ordre et les marque,
avec le numéro d'outbox comme clé d'idempotence afin que les répétitions après un plantage du relais soient écartées.
\iqlookfor{Outbox et clé d'idempotence.}
\end{solution}

\begin{solution}{iq:pl:messaging-and-event-driven-architecture:4}
Regrouper les messages par compte, de sorte que les messages d'un même compte aillent dans une même partition, et faire lire chaque partition par un seul consommateur du groupe
à la fois.
\iqlookfor{Clé vers partition~; un seul lecteur par partition.}
\end{solution}

\begin{solution}{iq:pl:messaging-and-event-driven-architecture:5}
Un \omterm{def:pl:messaging-and-event-driven-architecture:log}{journal partitionné} avec des topics par type d'événement, regroupés par compte ou par transaction~; des producteurs passant par des outboxes~; des consommateurs idempotents ou atomiques avec
leurs \omterm{def:pl:messaging-and-event-driven-architecture:log}{positions du consommateur}~; une rétention assortie d'instantanés pour le rejeu~; des \omterm{def:pl:messaging-and-event-driven-architecture:dlq}{files des messages en échec} surveillées~; le retard des consommateurs et les \omterm{def:pl:post-trade-systems:recon}{rapprochements} avec la drop copy comme
indicateurs de santé.
\iqlookfor{Journal, outbox, consommateurs idempotents, \omterm{def:pl:post-trade-systems:recon}{rapprochement}.}
\end{solution}
